\subsection{Additional known-parameter benchmark}

A separate fixed-parameter comparison evaluates six estimators on 32
final-validation data sets: four configurations with eight replications each,
for 192 valid fits. The methods receive $(r_0,g_0,\alpha_0)$; TKM and TCLUST,
which have no latent-rank parameter, receive $(g_0,\alpha_0)$. REG-HD-TFKM
uses HIER4 start arbitration and REG-HD-TRKM uses RV--CV. Each fit uses
100 starts. This experiment isolates fitting performance at known parameters;
it does not evaluate automatic structural selection on the 41-cell design.

\begin{table}[htbp]
\centering
\caption{Fixed-parameter comparison on 32 final-validation data sets
(four configurations, eight replications each). ARI is evaluated on all
 genuine observations, including those trimmed by the fit. Runtime is the
recorded mean seconds per fit.}
\label{tab:supp-competitors}
\small
\begin{tabular}{lrrr}
\toprule
Method & Mean ARI & Median ARI & Runtime\\
\midrule
TKM & .9782 & 1.0000 & .524\\
PCA--TKM & .9737 & 1.0000 & .101\\
ROBPCA--TKM & .9851 & 1.0000 & .129\\
TCLUST & .9779 & 1.0000 & 17.203\\
REG-HD-TFKM & .9836 & 1.0000 & 29.821\\
REG-HD-TRKM & .9834 & 1.0000 & 27.280\\
\bottomrule
\end{tabular}
\end{table}

Mean genuine-observation ARI ranges from $.974$ to $.985$, and every method
has median ARI one in this panel. These are descriptive results for a limited
fixed-parameter comparison, not evidence of universal dominance. In bridge
samples, supplying the nominal contamination fraction does not by itself
remove allocation or trimming ambiguity. The distinction between fitting
quality, structural selection and exact trimming recovery therefore remains
important even when partition agreement is high.


